% !TeX root = ../../hodge-classes.tex
\section{Purpose and notation}
This supporting note records explicit Hodge-algebraicity arguments for the
fixed field $M=\mathbb Q(\zeta_{42})$. The strongest statement is
Theorem~\ref{eight:thm:all}: arbitrary finite mixed products of the eight
rank-one $M$ CM-type orbits. The earlier mixed-power and six-orbit results
are retained as proof stages, rather than presented as separate novelty
claims. They supply source tensors used in the final argument.

A rank-one $M$ sixfold means a complex abelian sixfold whose rational first
cohomology has dimension one over $M$. Embeddings are labelled by units
modulo $42$. We use $K=\mathbb Q(\sqrt{-3})$, $L=M^+$ and
$\Phi_B=\{1,5,23,25,29,31\}$. The symbols $A_A,A_B,A_C$ in the packet
corollary refer to the geometric fibers of Paper I, Theorem~\externalref{packet}{thm:A}.
That application uses only their stated isogeny classes; the abstract
mixed-product theorem does not rely on the packet intersection search.

The route is: explicit algebraic input cycles, a saturated integral
character-relation lattice, tensor products and contractions, and rational
descent. A rational relation basis alone would not justify the integral
step. Computational relation checks supplement the arguments; they do
not by themselves prove cycle algebraicity.

OpenAI's proposed general CM theorem \cite{OpenAICM2026}, if valid,
implies the conclusions here. Andr\'e's reduction and the split-Weil
literature \cite{Andre1992,MilneHodge2020,Markman2025,Schoen1998}, and the
Fermat antecedents discussed by Saller \cite{Saller2026}, must be included
in assessing novelty. This note claims neither a new general CM range
nor an independent verification of the OpenAI theorem. The reported Aoki
route is not an adopted proof input. Its role is to document the explicit
fixed-field construction supporting the two arithmetic-geometric papers.

Theorem~\ref{ext:thm:mixed} is proved by the algebraic input construction
and the synchronized relation calculation of
Sections~\ref{ext:sec:inputs}--\ref{ext:sec:lattice}.
The individual-powers statement follows from the eight-type analysis
following that calculation; the last section proves the full mixed result.
